% John Dowland to Sir Robert Cecil, Nuremberg, 10 November 1595.
% Original text as transcribed at walterbitner.com (after Pinto); words the
% manuscript crosses out are shown with \strike. Modern English translation
% from john-dowland-letter-modern-english.docx.
%
% The letter is cut into pieces so each marked passage (\Oi ... \Ovii and
% \Ti ... \Tvii) can be used both in the full letter and on its own page.
% \OgX / \TgX are the unmarked stretches between passages.
% Passages in letter order: 1, 5, 2, 4, 3, 6, 7.

% ---------------------------------------------------------------- original

\newcommand{\OiA}{Ryght honorable: as I have bin most bounde unto your honor, so I most humblie desire your honor to pardon my boldnes, and mak my choic of your honor to let you understand my bounden duty and desire of Gods preservation of my most dear soveraigne Quene and contry: whom I besech God ever to bless and to confounde all their enemies what and whomsoever.}
\newcommand{\OiB}{Fiften yeares since I was in France Servant to Sir Henry Cobham who was Imbassador for the Quenes majestie, and lay in Parys, wher I fell aquainted with on Smith a preist, and on Morgan somtims of her majesties chapell, on Verstigan who brake out of Englande beinge aprehended, and on Moris a Welchman that was our porter, who is at Rome, these men thrust many idle toies into my hed of relygion, sainge that the papists was the truthe and ours in England all falce, and I being but yonge their faire words overecht me and I beleved with them.}

\newcommand{\OgA}{Within ij years after I cam into Englande wher I saw men of that faction condemde and executed, which I thought was great injustice taking relygion for the only cause, and when my best frends wold perswade me I wold not beleve them.}

\newcommand{\Ov}{Then in time passing on Mr.\ Johnson died, and I becam an humble sutor for his place (thincking my selfe most worthiest), wherin I found many goode and honorable frends that spake for me, but I saw that I was like to goe without it, and that any mygt have preferment but I. Wherby I began to sounde the cause, and gest that my relygion was my hinderance. Whearupon, my mynde beinge trobled, I desired to get beyond the seas,}

\newcommand{\OgB}{which I durst not attempt without lycence from som of the Privie Counsell, for fear of being taken and so my hav extreame punishment.}

\newcommand{\OiiA}{And accordinge as I desired ther cam a letter to me out of Germany from the Duke of Brunswicke, wherupon I spake to your honor and to my Lord of Essex, who willingly gav me both your hands (for which I wold be glad if ther wear any service in me that yor honors coulde comande).}
\newcommand{\OiiB}{When I cam to the Duke of Brunswicke he usde me kindly, and gave me a ritch chaine of golde, xxiij~£ in mony with velvet and saten and gold lace to make me aparell, with promyse that if I wolde serv him he wolde gev me as mutch as any prince in the worlde.}
\newcommand{\OiiC}{From thenc I went to the Landgrave of Hessen, (who gav me the gretest welcome that myght be for on of my qualyty) who sent a ringe into England to my wiff valued at xx~£ sterlinge, and gave me a great standing cupe with a cover gilt, full of dolers with many great offers for my service.}
\newcommand{\OiiD}{And from thence I had great desire to see Italy and cam to Venyce and from thence to Florence whear I plaid befor the Duke and got great favors.}

\newcommand{\Oiv}{And on evenynge I was walkinge upon the piazzo in Florence a jentle man tolde me that he espied a English preist and that his name was Skidmore and sonne and heir to Sir John Skidmore of the Courte, so I being intended to go unto Rome to study with a famous musicion named Luca Marenzio, stept to this Mr.\ Skidmor the preist and asked him if he was an Englysh man, and he told me yea: and whose sonne he was, and I telling him my name he was very glad to se me, so I tolde him I wold goe to Rome and desired his help for my saffty. For said I, “if they shold mystake ther, my fortune wear hard, for I hav bin thrust off of all good fortunes because I am a Catholicke at home.}

\newcommand{\Oiii}{For I hard that Her Majestie beinge spake to for me, saied, I was a man to serve any princ in the world, but I was an obstinat papist.” Whearunto he answered, “Mr.\ Dowlande, if it be not so, make her words true.” So in farther talke we spake of preists, and I told him that I did not thinck \strike{thinck} it true that any preists (as we said in England) wolde kill the Quen, or onc goe about to touche her finger, and said I whatsoever my Relygion be will neither medle nor mak with any thyng thear done, so that they do not any thing against the Quen. Whearunto he answered that I spake as a good Subiect to her majestie, but said he in Rome you shall Englyshmen your own Contrymen speak most hardly of her and wholy to overthrow her and all Englande, and those be the Jesuits, said he, who are of the Spanyshe faction: moreover said he we hav many Jares with them, and withall wished to God the Quen wear a Catholycke.}

\newcommand{\OgCa}{And said he, “to defend my Contry against the Spanyards I wold com into England and bear a pike on my sholders.”}
\newcommand{\OgCb}{Amonge our talk he tolde me that he had order to atache \strike{div} dyvers Inglish Jentlemen, and that he had bine iij years from Englande. So I brought him to his lodging dore, whear he told me that thear was ix prests com from Rome to go for Englande. He cam but the day before to Florenc, and I thinck they cam altogether.}
\newcommand{\OgCc}{He told me that he wold stay ther in the town and study in an abby caled Sancta Maria Nouella, and that he must kep in for on monethe, and that he wold writ letters of me to Rome, which I shold receve very shortly. But I heard not of him in a moneth after, and then there came two English friers to my lodging the on was an Englishman named Balye, a Yorksheirman.}
\newcommand{\OgCd}{The next day after my spech with Skydmore, I dyned with my Lord Gray and dyvers other jentlemen, whom I told of my speach with Skidmor, gevinge them warnynge. Wherupon my Lord Gray went to Siena, and the rest dispersed them selves.}
\newcommand{\OgCe}{More over I tolde my Lord Gray howsoever I was for relygion, if I did perceve anything in Rome that ether toucht her majestie of the state of England I wold gev nottice of it thoughe it wer the loss of my liff. Which he liked well and bade me kepe that secret.}
\newcommand{\OgCf}{This Frier Baylie before named delivered me a letter which I hav hear sent unto your honor, which letter I brake open befor Mr.\ \strike{bodly} Josias Bodly, and shewed what was written in it to him and divers other, after this, this frier Bayly tolde he had Receved letters From Rom to hasten me forward, and tolde me that my discontentment was knowne at Rom, and that I sholde hav a large pention of the Pope, and that His Holynes and all the Cardinales woulde make wonder full mutch of me.}
\newcommand{\OgCg}{Therupon I tolde him of my wiff and children how to get them to me, wherunto he tolde me that I shold have acquaintance with sutche as sholde bringe them over if she had any wilingnes, or elce they wolde lose ther lives. For ther cam those into Englande for sutch purposes. “For,” quothe he “Mr.\ Skidmor brought out of Englande at his last being ther xvij persons both men and women, for which the Bishope wepes when he sees him for Joy.”}

\newcommand{\OviA}{After my departure I caled to mynde our conference and got me by my self and wept hartely, to se my fortun so harde that I sholde becom sarvant to the gretest enemye of my prince, contry, wiff, children, and frends, for wante, and to make me like them selves.}
\newcommand{\OviB}{God He knoweth I never loved treason nor trechery nor never knew any, nor never heard any Mass in Englande, which I finde is great abuse of the peple for on my soule I understande it not, wherfor I have reformed my self to lyve acording to Her Majesties lawes as I was borne under Her Highnes, and that most humbly I doe crave pardon}

\newcommand{\OgDa}{, protestinge if ther wer any abylitie in me, I wold be most redy to make amends.}
\newcommand{\OgDb}{At Bolona I met with ij men the on named pierce an Iryshman, the other named Dracot, they are gon both to Rome, In Venice I hard an Italyan say, that he marveled that king Phillip had never a good frend in Inglande that with his dager wolde dispatch the Quens Majestie, but said he, God suffers her, in the end to gev her the greter overthrowe.}
\newcommand{\OgDc}{Right honorable, this hav I written that Her Majestie may knowe the vilany of thes most wicked prests and Jesuits, and to bewar of them. I thanck God I hav both forsaken them and their relygion which tendeth to nothing but distruction.}
\newcommand{\OgDd}{Thus I beseche God nyght and day to blesse and defend the Quens Majestie, and to confonde all her enemies and to preserve your honor and all the rest of her majesties most honorable Privie Counsell.}
\newcommand{\OgDe}{I thinck that Skidmore and the other preists ar all in England, for he staid not at Florenc as he said he wolde to me, and Frier Baylie tolde me that he was gon into France to study the lawe.}
\newcommand{\OgDf}{At Venyce and all along as I com in Germany men say that the Kinge of Spain is making gret preparation to com for England this next somer, wher}

\newcommand{\Ovii}{if it pleasde your honor to advise me by my poore wyff I wolde most wilingly lose my lyffe against them, most humbly besechinge your honor to pardon my ill writinge, and worse inditinge, and to thinke that I desir to serve my contry, and hope to hear of your good opinyon of me.}
\newcommand{\OviiDate}{From Nurnberge, this xth of November 1595}
\newcommand{\OviiClose}{Your honors most bounden}
\newcommand{\OSign}{for ever\\John Doulande}

% ------------------------------------------------------------- translation

\newcommand{\TiA}{Right Honourable Sir: Because I am deeply indebted to Your Honour, I most humbly ask you to forgive my presumption, and I choose Your Honour as the person to whom I may explain my solemn obligation and my desire that God preserve my most dear sovereign, the Queen, and my country.\fn{1} I pray that God will always bless them and confound every one of their enemies, whoever they may be.\fn{2}}
\newcommand{\TiB}{Fifteen years ago I was in France, in the service of Sir Henry Cobham, who was Her Majesty the Queen's ambassador and was residing in Paris.\fn{3} There I became acquainted with a priest named Smith; with Morgan, formerly a member of Her Majesty's Chapel; with Verstegan, who escaped from England after being apprehended; and with Morris, a Welshman who was our porter and who is now in Rome.\fn{4} These men filled my head with many foolish ideas about religion, saying that the papists possessed the truth and that our religion in England was entirely false. Because I was still young, their persuasive words deceived me, and I came to share their beliefs.\fn{5}}

\newcommand{\TgA}{Within two years after that I returned to England, where I saw men of that party condemned and executed. I believed this to be a great injustice, because I thought religion was the only reason for their punishment; and when my closest friends tried to persuade me otherwise, I would not believe them.\fn{6}}

\newcommand{\Tv}{As time passed, Mr Johnson died, and I humbly applied for his position, thinking myself the most worthy candidate.\fn{7} I found many good and honourable friends who spoke on my behalf, but I saw that I was likely to go without the position and that anyone might receive advancement except me. I therefore began to investigate the cause and guessed that my religion was the obstacle.\fn{8}

My mind being troubled by this, I wished to go abroad.}

\newcommand{\TgB}{I did not dare attempt it without permission from members of the Privy Council, for fear that I might be arrested and severely punished.\fn{9}}

\newcommand{\TiiA}{Just as I wished, a letter then came to me from Germany from the Duke of Brunswick. I therefore spoke to Your Honour and to my Lord of Essex, and you both willingly gave me your signatures and authorization.\fn{10} For this kindness I would be glad if there were any service within my power that Your Honours could command me to perform.}
\newcommand{\TiiB}{When I came to the Duke of Brunswick, he treated me kindly. He gave me a rich gold chain, twenty-three pounds in money, and velvet, satin and gold lace from which to make clothing. He also promised that, if I would serve him, he would give me as much as any prince in the world would give me.\fn{11}}
\newcommand{\TiiC}{From there I went to the Landgrave of Hesse, who gave me as great a welcome as could be given to a man of my profession and standing.\fn{12} He sent my wife in England a ring valued at twenty pounds sterling and gave me a large standing cup with a gilded cover, filled with thalers, together with many generous offers for my service.\fn{13}}
\newcommand{\TiiD}{From there I felt a great desire to see Italy. I came to Venice and then to Florence, where I played before the Duke and received great favour.\fn{14}}

\newcommand{\Tiv}{One evening, while I was walking in the piazza in Florence, a gentleman told me that he had spotted an English priest. He said that the priest's name was Skidmore and that he was the son and heir of Sir John Skidmore, a man of the royal court.\fn{15}

Because I intended to go to Rome to study with a famous musician named Luca Marenzio, I went up to Mr Skidmore, the priest, and asked whether he was an Englishman.\fn{16} He said that he was, and told me whose son he was. When I gave him my name, he was very glad to see me. I told him that I wished to go to Rome and asked for his help in securing my safety there.

I said, “If they should mistake me for something I am not when I arrive, my fortune would be hard indeed, because at home I have been excluded from every good opportunity on account of being considered a Catholic.\fn{17}}

\newcommand{\Tiii}{I have heard that, when Her Majesty was approached on my behalf, she said that I was a man capable of serving any prince in the world, but that I was an obstinate papist.”\fn{18}

He answered, “Mr Dowland, if that is not true, then make her words true.”\fn{19}

As our conversation continued, we spoke about priests. I told him that I did not believe it was true that any priests, as people claimed in England, would kill the Queen or even attempt to touch one of her fingers. I added, “Whatever my religion may be, I will neither involve myself nor take part in anything done there, provided that they do nothing against the Queen.”\fn{20}

He replied that I spoke like a loyal subject of Her Majesty. “But in Rome,” he said, “you will hear Englishmen, your own countrymen, speak most harshly of her and seek entirely to overthrow both her and all England. Those men are the Jesuits, who belong to the Spanish faction.” He added that his own group had many quarrels with them. At the same time, he wished to God that the Queen were a Catholic.}

\newcommand{\TgCa}{He then said, “To defend my country against the Spaniards, I would return to England and carry a pike on my shoulders.”\fn{21}}
\newcommand{\TgCb}{During our conversation, he told me that he had orders to apprehend various English gentlemen and that he had been away from England for three years.\fn{22} I accompanied him to the door of his lodging, where he told me that nine priests had come from Rome on their way to England. He himself had arrived in Florence only the previous day, and I think that they had all arrived together.}
\newcommand{\TgCc}{He told me that he would remain in the town and study in an abbey called Santa Maria Novella; that he must stay secluded there for one month; and that he would write letters about me, or on my behalf, to Rome, which I should receive very shortly.\fn{23} But I heard nothing from him for a month. Then two friars came to my lodging. One of them was an Englishman named Baylie, a Yorkshireman.\fn{24}}
\newcommand{\TgCd}{The day after my conversation with Skidmore, I dined with Lord Grey and several other gentlemen. I told them what Skidmore and I had discussed and warned them. Lord Grey consequently went to Siena, and the others dispersed.\fn{25}}
\newcommand{\TgCe}{I also told Lord Grey that, whatever my religion might be, if I discovered anything in Rome that affected Her Majesty or the English state, I would report it even if doing so cost me my life. He approved of this and instructed me to keep it secret.}
\newcommand{\TgCf}{The Friar Baylie mentioned above delivered a letter to me, which I have enclosed and sent to Your Honour. I broke the letter open in the presence of Mr Josias Bodley and showed him and several others what was written in it.\fn{26} After this, Friar Baylie told me that he had received letters from Rome urging me to proceed quickly. He said that my dissatisfaction was known in Rome, that I would receive a large pension from the Pope, and that His Holiness and all the cardinals would make a great deal of me.\fn{27}}
\newcommand{\TgCg}{I then asked him about my wife and children and how they could be brought to me. He told me that I would be introduced to people who could bring them over if my wife were willing; otherwise, those people would lose their lives, for men came into England for such purposes.\fn{28} “For,” he said, “when Mr Skidmore was last in England, he brought seventeen people out of the country, both men and women. The bishop weeps for joy whenever he sees him because of it.”\fn{29}}

\newcommand{\TviA}{After I left him, I recalled our conversation, withdrew by myself and wept bitterly. I grieved that my fortune should be so hard that poverty might force me to become the servant of the greatest enemy of my sovereign, my country, my wife, my children and my friends, and that they should try to make me like themselves.\fn{30}}
\newcommand{\TviB}{God knows that I have never loved treason or treachery, nor have I ever known of any. I have never attended Mass in England, which I now find to be a great deception of the people; upon my soul, I do not understand it.\fn{31} Therefore, I have reformed my conduct and resolved to live according to Her Majesty's laws, since I was born under Her Highness.\fn{32} For my past conduct I most humbly ask pardon,}

\newcommand{\TgDa}{declaring that, if I possess any ability to do so, I will be entirely ready to make amends.}
\newcommand{\TgDb}{At Bologna I met two men, one named Pierce, an Irishman, and the other named Dracot. They have both gone to Rome.\fn{33} In Venice I heard an Italian say that he was astonished King Philip had no good supporter in England who would dispatch Her Majesty the Queen with a dagger. He then said that God was allowing her to survive so that, in the end, He might give her an even greater overthrow.\fn{34}}
\newcommand{\TgDc}{Right Honourable Sir, I have written this so that Her Majesty may know the villainy of these most wicked priests and Jesuits and may be on her guard against them. I thank God that I have forsaken both them and their religion, which leads to nothing but destruction.\fn{35}}
\newcommand{\TgDd}{Thus I pray to God night and day to bless and defend Her Majesty the Queen, to confound all her enemies, and to preserve Your Honour and all the other members of Her Majesty's most honourable Privy Council.}
\newcommand{\TgDe}{I think that Skidmore and the other priests are now all in England, because he did not remain in Florence as he had told me he would. Friar Baylie told me that Skidmore had gone to France to study law, but I remain suspicious.\fn{36}}
\newcommand{\TgDf}{In Venice, and everywhere along my journey through Germany, people say that the King of Spain is making great preparations to come against England next summer.}

\newcommand{\Tvii}{If it pleased Your Honour to send me instructions through my poor wife, I would most willingly lose my life fighting against them.\fn{37}

I most humbly ask Your Honour to pardon my poor handwriting and my still poorer composition, and to believe that I desire to serve my country. I hope to hear that you hold a good opinion of me.\fn{38}}
\newcommand{\TviiDate}{From Nuremberg, this 10th day of November 1595}
\newcommand{\TviiClose}{Your Honour's most deeply obliged servant}
\newcommand{\TSign}{forever\\John Dowland}
